\section{光滑态射}\label{sec:ag-3-4}
% 原PDF p58.
前面已回顾过：代数闭域上的簇，若其所有点都正则，就称为光滑簇. 这里的目的是引入光滑性的相对版本. 可以给出若干种定义，我们将验证它们彼此等价.
第一种定义受微分几何启发.
\begin{definition}\label{item:ag-3-4-1}
设 $f:X\longrightarrow S$ 是态射，$x\in X$. 若存在 $x$ 的开邻域 $U$、整数 $n$ 及图
\begin{align*}
\begin{tikzcd}[ampersand replacement=\&]
U\arrow[r,"\varphi"]\arrow[dr,"f|_U"'] \& S\times\mathbb A^n\arrow[d]\\
\& S
\end{tikzcd}
\end{align*}
其中 $\varphi$ 是平展态射，则称 $f$ \emph{在 $x$ 处光滑}. 若 $f$ 在 $X$ 的所有点处都光滑，则称 $f$ 为\newterm{ag-smooth-morphism}{光滑态射}{smooth morphism}.
\end{definition}
与微分几何的类比可能使人预期的要求相比，这里允许 $\varphi$ 为平展态射，而不要求它为开浸入. 首先要做的是，当 $S$ 为代数闭域的谱时，将此定义与前面的定义作比较.
\begin{proposition}\label{item:ag-3-4-2}
设 $k$ 是代数闭域，$X$ 是有限型 $k$-概形，$x$ 是 $X$ 的闭点. 则结构态射 $f:X\longrightarrow\operatorname{Spec}(k)$ 在点 $x$ 处光滑，当且仅当局部环 $\mathcal O_{X,x}$ 正则.
\end{proposition}
\begin{proof}
正向蕴含来自以下观察：对 $\mathbb A_k^n$ 的每一点 $y$，环 $\mathcal O_{\mathbb A_k^n,y}$ 都正则；而正则性在平展态射下保持，参见命题~\ref{item:ag-3-3-13} ii).
证明反向蕴含. 设 $\mathcal O_{X,x}$ 正则，选取 $x$ 的开邻域 $U$，以及在 $x$ 处为零、且在 $\mathfrak m_x/\mathfrak m_x^2$ 中的像构成一组基的函数 $f_1,\ldots,f_n$（于是 $n=\dim_x(X)$）. 这些函数定义态射 $U\xrightarrow\varphi\mathbb A_k^n$. 由``正则''的定义，局部态射 $\mathcal O_{\mathbb A_k^n,\varphi(x)}\longrightarrow\mathcal O_{X,x}$ 诱导分次环同构 $\operatorname{gr}^{\mathfrak m_{\varphi(x)}}\mathcal O_{\mathbb A_k^n}\xrightarrow\sim\operatorname{gr}^{\mathfrak m_x}\mathcal O_{X,x}$. 由于剩余域扩张平凡（事实上，由零点定理，$\kappa(x)=\kappa(y)=k$），利用~\ref{item:ag-3-3-13}(i) 可知，$\varphi$ 在 $x$ 处平展. 最后只须适当缩小 $U$.
\end{proof}
\begin{remark}\label{item:ag-3-4-3}
\begin{enumerate}[label=(\roman*)]
\item 若 $k$ 不是代数闭域，上述证明仍表明：若 $f$ 在 $x$ 处光滑，则 $\mathcal O_{X,x}$ 正则. 只有当 $\kappa(x)$ 是 $k$ 的可分扩张时，逆命题才成立. 特别地，若 $k$ 为完美域，则对有限型 $k$-概形，``在 $k$ 上光滑''等价于``正则''.
\item 在代数闭域上，$X$（姑且设其连通）在 $k$ 上光滑，当且仅当 $\Omega_{X/k}$ 自由且秩为 $\dim(X)$. 以后将看到，这一等价关系在任意域上仍成立. 特别地，$X$ 在域 $k$ 上光滑，当且仅当 $X_{\bar k}$ 在 $\bar k$ 上光滑.\unskip\nobreak\hfill\ensuremath{\diamond}
\end{enumerate}
\renewcommand{\envenddiamond}{}
\end{remark}
% 原PDF p59.
\begin{remark}\label{item:ag-3-4-4}
\begin{enumerate}[label=(\roman*)]
\item 每个平展态射都光滑.
\item 光滑态射在复合与基变换下保持.\unskip\nobreak\hfill\ensuremath{\diamond}
\end{enumerate}
\renewcommand{\envenddiamond}{}
\end{remark}
\begin{proposition}[逐纤维判别准则]\label{item:ag-3-4-5}
设 $f:Y\longrightarrow X$ 是 $S$-概形的态射，$y$ 是 $Y$ 的一点，$x$ 是其在 $X$ 中的像，$s$ 是其在 $S$ 中的像. 则有：
\begin{enumerate}[label=\roman*)]
\item $f$ 在点 $y$ 处非分歧，当且仅当 $f\times\kappa(s):Y_s\longrightarrow X_s$ 在点 $y$ 处非分歧.
\item 再假设 $X$、$Y$ 在 $S$ 上平坦. 则 $f$ 在 $y$ 处平坦（分别为平展、光滑），当且仅当 $f\times\kappa(s):Y_s\longrightarrow X_s$ 在 $y$ 处平坦（分别为平展、光滑）.
\end{enumerate}
\end{proposition}
\begin{proof}
\textup{(i)} 有 $\mathcal O_{Y_s,y}=\mathcal O_{Y,y}/\mathfrak m_s\mathcal O_{Y,y}$，$\mathcal O_{X_s,x}$ 也同样. 因而 $\mathcal O_{Y,y}/\mathfrak m_x\mathcal O_{Y,y}=\mathcal O_{Y_s,y}/\mathfrak m_{X_s,x}\mathcal O_{Y,y}$. 判别准则~\ref{item:ag-3-3-4} ii) 于是说明，$g$ 在 $y$ 处非分歧当且仅当 $g\times\kappa(s)$ 在该点非分歧.
\textup{(ii)} 对这三个性质，``仅当''方向均来自它们在基变换下保持这一事实，故只须证明``当''方向.
先处理平坦性. 多次应用判别准则~\ref{item:ag-3-3-10} ii). 事实上，$\mathcal O_{Y,y}$ 与 $\mathcal O_{X,x}$ 在 $\mathcal O_{S,s}$ 上的平坦性给出下列同构：
\begin{align*}
\mathcal O_{X,x}\otimes_{\mathcal O_{S,s}}\operatorname{gr}^{\mathfrak m_s}(\mathcal O_{S,s})&\xrightarrow\sim\operatorname{gr}^{\mathfrak m_s}(\mathcal O_{X,x}),\\
\mathcal O_{Y,y}\otimes_{\mathcal O_{S,s}}\operatorname{gr}^{\mathfrak m_s}(\mathcal O_{S,s})&\xrightarrow\sim\operatorname{gr}^{\mathfrak m_s}(\mathcal O_{Y,y}),
\end{align*}
由此可知，典范态射
\begin{align*}
\mathcal O_{Y,y}\otimes_{\mathcal O_{X,x}}\operatorname{gr}^{\mathfrak m_s}(\mathcal O_{X,x})\xrightarrow\sim\operatorname{gr}^{\mathfrak m_s}(\mathcal O_{Y,y})
\end{align*}
是同构. 再结合 $\mathcal O_{Y,y}/\mathfrak m_s$ 在 $\mathcal O_{X,x}/\mathfrak m_s$ 上平坦的假设即可得出结论. 具体取 $A=\mathcal O_{X,x}$、$M=\mathcal O_{Y,y}$、$I=\mathfrak m_s\mathcal O_{X,x}$.
平展性由已经处理过的平坦性和非分歧性得到.
接着处理光滑情形. 图
\begin{align*}
\begin{tikzcd}[ampersand replacement=\&]
U_s\arrow[r,"\varphi_s"]\arrow[d,hook] \& \mathbb A_{X_s}^n\arrow[d]\\
Y_s\arrow[r] \& X_s
\end{tikzcd}
\end{align*}
由 $n$ 个函数 $f_i\in\Gamma(U_s,\mathcal O_{U_s})$ 给出. 仍记 $f_i$ 为这些函数在 $\mathcal O_{Y_s,y}$ 中的芽，并在 $\mathcal O_{Y,y}$ 中选取它们的提升 $\widetilde f_i$. 这些芽延拓到 $y$ 的某个邻域 $U$，并定义图
\begin{align*}
\begin{tikzcd}[ampersand replacement=\&]
U\arrow[r,"\varphi"]\arrow[d,hook] \& \mathbb A_X^n\arrow[d]\\
Y\arrow[r] \& X
\end{tikzcd}
\end{align*}
由已证结论，因 $\varphi_s$ 在 $y$ 处平展，$\varphi$ 也在 $y$ 处平展. 必要时缩小 $U$ 后，就得到光滑定义中的图.
\end{proof}
% 原PDF p60.
\begin{remark}\label{item:ag-3-4-6}
特别取 $S=X$，便知态射 $f:Y\longrightarrow X$ 光滑当且仅当它平坦且纤维光滑.\unskip\nobreak\hfill\ensuremath{\diamond}
\renewcommand{\envenddiamond}{}
\end{remark}
\begin{proposition}[局部性质]\label{item:ag-3-4-7}
若 $f:Y\longrightarrow X$ 在 $y$ 处光滑，则 $\mathcal O_{Y,y}$ 正则（分别为 Cohen--Macaulay、正规、既约），当且仅当 $\mathcal O_{X,x}$ 具有相应性质.
\end{proposition}
\begin{proof}
这来自这些断言对平展态射及态射 $\mathbb A_X^n\longrightarrow X$ 均成立.
\end{proof}
\begin{proposition}\label{item:ag-3-4-8}
若 $f:X\longrightarrow S$ 光滑，则 $f$ 微分光滑. 特别地，$\Omega_{X/S}$ 局部自由，其秩为 $f$ 的相对维数.
\end{proposition}
\begin{proof}
这来自平展态射及态射 $\mathbb A_S^n\longrightarrow S$ 都微分光滑，以及平坦且微分光滑的态射的复合仍平坦且微分光滑. 对于最后一点，只须使用命题~\ref{item:ag-3-2-9}证明中的图.
\end{proof}
现在要考察：光滑概形的闭子概形何时光滑.
\begin{theorem}\label{item:ag-3-4-9}
设 $X$ 是光滑 $S$-概形，$Y$ 是 $X$ 的闭子概形，其理想为 $\mathcal I$，且 $x\in X$. 则下列条件等价：
\begin{enumerate}[label=\roman*)]
\item $Y$ 在 $x$ 处光滑.
\item 存在 $S$-概形的图
\begin{align*}
\begin{tikzcd}[ampersand replacement=\&]
Y\arrow[d,hook] \& U\cap Y\arrow[l,hook]\arrow[r,"\varphi|_{U\cap Y}"]\arrow[d,hook] \& \mathbb A_S^{n-p}\arrow[d,hook]\\
X \& U\arrow[l,hook]\arrow[r,"\varphi"'] \& \mathbb A_S^n
\end{tikzcd}
\end{align*}
其中 $U$ 是 $X$ 的开集，$\varphi$ 是平展态射，且两个方块均为笛卡尔方块.（此时 $p=\operatorname{Codim}(Y,X)$.）
\item \emph{Jacobian 判别准则.} 存在 $\mathcal I_x$ 作为 $\mathcal O_{X,x}$-模的一组生成元 $g_1,\ldots,g_p$，使 $dg_i$ 在 $\Omega_{X/S,x}\otimes\kappa(x)=\Omega_{\mathcal O_{X,x}/\mathcal O_{S,s}}\otimes\kappa(x)$ 中线性无关.（回顾，$d:\mathcal O_{X,x}\longrightarrow\Omega_{\mathcal O_{X,x}/\mathcal O_{S,s}}$ 是万有导子.）
\item 列
\begin{align*}
0\longrightarrow\mathcal N_{Y/X}=\mathcal I/\mathcal I^2\xrightarrow d i^*\Omega_{X/S}\longrightarrow\Omega_{Y/S}\longrightarrow0
\end{align*}
在 $x$ 处正合，且由局部自由模组成.
\end{enumerate}
\end{theorem}
\begin{proof}
对 $Y=\mathbb A_S^p$ 到 $X=\mathbb A_S^n$ 的浸入，性质 i)、ii) 和 iv) 均成立. 由于平展态射在 $\Omega$ 上诱导同构，参见注~\ref{item:ag-3-3-8} v)，可知 ii) 蕴含这三个性质.
证明 iii)$\Rightarrow$ii). 可在 $\mathcal O_{X,x}$ 中找到函数 $g_{p+1},\ldots,g_n$，使它们的微分将线性无关族 $dg_i$（$i\le p$）补成 $\Omega_{X/S}$ 在 $x$ 处的一组基. 这些芽
% 原PDF p61.
延拓为 $x$ 的某个邻域上的函数，并定义 ii) 中那样的图. 由构造，$\varphi$ 在点 $x$ 的邻域上诱导同构 $\varphi^*\Omega_{\mathbb A_S^n/S}\xrightarrow\sim\Omega_{X/S}$，从而在此邻域内有 $\Omega_{U/\mathbb A_S^n}=0$. 还须说明 $\varphi$ 在 $x$ 处平坦. 由逐纤维判别准则，归结到 $S$ 为域的谱的情形. 此时已知 $\mathcal O_{X,x}$ 为正则环，$\mathcal O_{\mathbb A^n,\varphi(x)}$ 也同样正则. 由构造，态射 $\varphi$ 诱导分次环同构
\begin{align*}
\kappa(x)\otimes\operatorname{gr}^{\mathfrak m_{\varphi(x)}}(\mathcal O_{\mathbb A_S^n,\varphi(x)})\xrightarrow\sim\operatorname{gr}^{\mathfrak m_x}(\mathcal O_{X,x}).
\end{align*}
由平坦性判别准则~\ref{item:ag-3-3-10}得出结论.
证明 i)$\Rightarrow$iii). 在正合列
\begin{align*}
\mathcal I/\mathcal I^2\xrightarrow d i^*\Omega_{X/S}\longrightarrow\Omega_{Y/S}\longrightarrow0
\end{align*}
中，后两个模在 $x$ 处自由. 因而可在 $\mathcal I_x$ 中选取芽 $g_1,\ldots,g_p$，使它们在 $d$ 下的像构成 $i^*\Omega_{X/S}\longrightarrow\Omega_{Y/S}$ 的核的一组基. 这些芽延拓为 $x$ 的某个邻域上的函数，由已经证明的蕴含，它们定义 $X$ 的一个 $S$-光滑闭子概形 $Y'$. 在 $x$ 的邻域内，由构造有 $Y\subset Y'$，因而由相对维数相等可得 $Y=Y'$.
最后，蕴含 iv)$\Rightarrow$iii) 是显然的.
\end{proof}
\begin{corollary}\label{item:ag-3-4-10}
设 $f:X\longrightarrow S$，$f':X'\longrightarrow S'$ 是它沿平坦 $S$-概形 $S'$ 的基变换. 设 $x'\in X'$ 的像为 $x\in X$. 则 $f$ 在 $x$ 处光滑，当且仅当 $f'$ 在 $x'$ 处光滑.
\end{corollary}
\begin{proof}
可设 $X$ 仿射，并在仿射概形 $S$ 上有限型. 将 $X$ 实现为 $\mathbb A_X^N$ 的闭子概形，再使用判别准则 (iv).
\end{proof}
\begin{remark}\label{item:ag-3-4-11}
\begin{enumerate}[label=(\roman*)]
\item 因此，可在忠实平坦基变换后检验光滑性.
\item 此推论说明了注~\ref{item:ag-3-4-3}中的断言.\unskip\nobreak\hfill\ensuremath{\diamond}
\end{enumerate}
\renewcommand{\envenddiamond}{}
\end{remark}
下面给出定理~\ref{item:ag-3-1-10}的 iii) 的相对版本.
\begin{theorem}\label{item:ag-3-4-12}
设 $X$ 是 $S$-概形，$Y$ 是 $X$ 的闭子概形，其理想为 $\mathcal I$，且 $x\in X$. 假设 $Y$ 在 $x$ 处光滑. 则以下条件等价：
\begin{enumerate}[label=\roman*)]
\item $X$ 在 $x$ 处光滑.
\item $Y$ 到 $X$ 的浸入正则.
\end{enumerate}
\end{theorem}
\begin{proof}
蕴含 i)$\Rightarrow$ii) 来自前一定理的判别准则 ii)，以及浸入 $\mathbb A_S^p\hookrightarrow\mathbb A_S^n$ 正则这一事实.
证明逆向时，在局部工作. 将 $\mathcal I$ 的正则生成元组 $g_1,\ldots,g_p$ 补以 $X$ 上的函数 $g_{p+1},\ldots,g_n$，使它们在 $Y$ 上的限制（局部地）定义平展态射 $Y\longrightarrow\mathbb A_S^{n-p}$. 这样（局部地）得到态射
% 原PDF p62.
$X\longrightarrow\mathbb A_S^n$，它嵌入笛卡尔图
\begin{align*}
\begin{tikzcd}[ampersand replacement=\&]
Y\arrow[d,hook] \& U\cap Y\arrow[l,hook]\arrow[r,"\varphi|_{U\cap Y}"]\arrow[d,hook] \& \mathbb A_S^{n-p}\arrow[d,hook]\\
X \& U\arrow[l,hook]\arrow[r,"\varphi"'] \& \mathbb A_S^n
\end{tikzcd}
\end{align*}
还须证明 $\varphi$ 平展. 因 $\varphi|_{Y\cap U}$ 平展，有 $(\Omega_\varphi)_x=(i^*\Omega_\varphi)_x=(\Omega_{\varphi|_{Y\cap U}})_x=0$，可见 $\varphi$ 非分歧. 对于平坦性，使用判别准则~\ref{item:ag-3-3-10} ii)，取 $A=\mathcal O_{\mathbb A_S^n,\varphi(x)}$、$B=M=\mathcal O_{X,x}$，并取理想 $I$ 使 $\mathcal O_{\mathbb A_S^{n-p},\varphi(x)}=A/I$. 于是 $B/IB=M/IM=\mathcal O_{Y,x}$ 确实在 $A/I$ 上平坦，而 $IB$ 正则的假设保证态射 $B\otimes_A\operatorname{gr}^{I}(A)\longrightarrow\operatorname{gr}^{IB}(B)$ 是同构.
\end{proof}
\begin{corollary}\label{item:ag-3-4-13}
设 $X$ 是概形，$Y$ 是闭子概形. 假设 $X$ 与 $Y$ 在 $S$ 上光滑. 则爆破 $\widetilde X_Y$ 及其例外除子 $Y'$ 在 $S$ 上光滑.
\end{corollary}
\begin{proof}
由前一定理，$Y$ 正则浸入 $X$. 根据命题~\ref{item:ag-3-1-7}，例外除子 $Y'$ 等同于 $\mathbb P(\mathcal N_{Y/X})$，故在 $Y$ 上光滑，并由复合可知在 $S$ 上光滑. 又因它是除子，故正则浸入 $\widetilde X_Y$；由前一定理，后者在 $S$ 上光滑.
\end{proof}
\begin{corollary}\label{item:ag-3-4-14}
$f:X\longrightarrow S$ 在 $x$ 处光滑，当且仅当它在 $x$ 处平坦且微分光滑.
\end{corollary}
\begin{proof}
已经证明了正向蕴含. 现设 $f$ 平坦且微分光滑. 由上面的推论，只须证明，$f$ 在任意忠实平坦扩张后成为光滑态射. 因而可设存在 $f$ 的截面 $i:S\hookrightarrow X$ 及一点 $s\in S$，使 $x=i(s)$（沿 $S'=X$ 作基变换）. 以下笛卡尔图
\begin{align*}
\begin{tikzcd}[ampersand replacement=\&]
S\arrow[r,hook,"i"]\arrow[d,"i"'] \& X\arrow[r,"f"]\arrow[d,"\operatorname{Id}\times i\circ f"] \& S\arrow[d,"i"]\\
X\arrow[r,hook,"\Delta_f"'] \& X\times_SX\arrow[r,"\pi_2"'] \& X
\end{tikzcd}
\end{align*}
以及命题~\ref{item:ag-3-1-14}表明，浸入 $S\hookrightarrow X$ 在 $x$ 处正则. 由于 $S$ 显然在 $S$ 上光滑，由前面的结果可知，$X$ 在 $x$ 处在 $S$ 上光滑.
\end{proof}\begin{numpara}[光滑性与形式光滑性]\label{item:ag-3-4-15}
Jacobian 判别准则原则上适于判定由方程给出的概形是否光滑. 然而，许多有趣的概形并非由方程给出，或至少用方程并不是理解它们的最佳方式. 模空间（如丛、Abel 簇、曲线等的模空间）就是这样的例子. 对于这些情形，有另一种更适合的光滑性判别准则.
% 原PDF p63.
若对任意仿射概形 $T$、任意由幂零理想定义的子概形 $T_0$，以及任意交换方块
\begin{align}
\begin{tikzcd}[ampersand replacement=\&]
T_0\arrow[r,"u_0"]\arrow[d] \& X\arrow[d,"f"]\\
T\arrow[r,"v"']\arrow[ur,dotted,"u"] \& S
\end{tikzcd}\tag{3.4.15.1}\label{eq:ag-3-4-15-1}
\end{align}
存在一个（分别为至多一个、恰好一个）使两个三角形交换的对角箭头 $u$，则称态射 $f:X\longrightarrow S$ \newterm{ag-formally-smooth-morphism}{形式光滑}{formally smooth morphism}（分别为\newterm{ag-formally-unramified-morphism}{形式非分歧}{formally unramified morphism}、\newterm{ag-formally-etale-morphism}{形式平展}{formally étale morphism}）. 称 $u$ 为``提升''. 注意：
\begin{itemize}
\item 对于``形式平展''和``形式非分歧''，可去掉 $T$ 仿射的假设，因为唯一性使提升能够粘合.
\item 通过显然的逐次分解，可设定义 $T_0$ 在 $T$ 中的理想 $\mathcal I$ 的平方为零，即 $\mathcal I^2=0$.
\end{itemize}
\end{numpara}
\begin{theorem}\label{item:ag-3-4-16}
态射 $f$ 光滑（分别为平展、非分歧），当且仅当它形式光滑（分别为形式平展、形式非分歧）.（回顾，已假设 $f$ 有限型.）
\end{theorem}
\begin{proof}
先看仿射情形，即考察环的图
\begin{align*}
\begin{tikzcd}[ampersand replacement=\&]
B\arrow[r,"u_0"]\arrow[dr,dotted,"u"] \& R/I\\
A\arrow[u,"f"]\arrow[r,"v"'] \& R\arrow[u]
\end{tikzcd}
\end{align*}
其中 $I$ 为平方为零的理想. 注意，$I=I/I^2$ 自然是 $R/I$-模，从而是 $B$-模. 设 $u,u'$ 是图中 $u_0$ 的两个提升. 计算可知，差 $u-u':B\longrightarrow I$ 是 $B$ 到 $I$ 的 $A$-导子. 因而
\begin{align}
\begin{aligned}
(u-u')\in\operatorname{Der}_A(B,I)&=\operatorname{Hom}_B(\Omega_{B/A},I)\\
&=\operatorname{Hom}_{R/I}(R/I\otimes_B\Omega_{B/A},I).
\end{aligned}\tag{3.4.16.2}\label{eq:ag-3-4-16-2}
\end{align}
\emph{非分歧情形.} 将上述计算应用于 $R=(B\otimes_AB)/I_{B/A}^2$ 和 $I=I_{B/A}/I_{B/A}^2=\Omega_{B/A}$，并取由 $b\in B\mapsto b\otimes1\in B\otimes B$ 给出的提升 $u$，就得到仿射情形下非分歧与形式非分歧的等价性. 在一般情形，只须用仿射开集之间的态射覆盖 $f$，即可得出结论，因为一个可能存在的提升若处处局部唯一，显然也全局唯一.
\emph{平展情形.} 证明平展态射形式平展. 从图~\eqref{eq:ag-3-4-15-1}出发，令 $X_T=X\times_ST$、$X_{T_0}=X\times_ST_0$. 已知平展态射 $X_{T_0}\xrightarrow{f_{T_0}}T_0$ 的一个截面 $u_0$，须找到平展态射 $X_T\xrightarrow{f_T}T$ 的一个截面 $u$. 注意，截面 $u_0$ 既是闭浸入又是开浸入，因为它是对角态射 $\Delta_{f_{T_0}}$ 沿态射 $\operatorname{Id}\times(u_0\circ f):$
% 原PDF p64.
$X_{T_0}\longrightarrow X_{T_0}\times_{T_0}X_{T_0}$ 的基变换. 因而，$u_0$ 是 $T_0$ 到 $X_{T_0}$ 的一个开闭子集 $Y_0$ 的同构. 所以 $f_{T_0}$ 诱导同构 $Y_0\xrightarrow\sim T_0$. 取 $X_T$ 的开闭子概形 $Y$，使其底层空间与 $Y_0$ 相同. 只须证明 $f_T$ 诱导同构 $Y\xrightarrow\sim T$. 而态射 $f_T$ 是万有同胚（即任意基变换后仍是同胚），因为 $f_{T_0}$ 是万有同胚. 特别地，$f_T$ 紧合，又因其拟有限，故为有限态射. 此外，$f_T$ 的剩余域扩张平凡，因为这些扩张与 $f_{T_0}$ 的相同. 因而，对 $t\in T$ 及其逆像 $y\in Y$，局部同态 $\mathcal O_{T,t}\longrightarrow\mathcal O_{Y,y}$ 有限，并在完备化上诱导同构. 所以它是同构，$f$ 确实诱导同构 $Y\xrightarrow\sim T$.
为帮助直观理解，可考虑以下特例：设 $S$ 仿射，其环为 $A$，$X$ 仿射，其环 $B$ 形如 $A[t]=A[T]/P(T)$，其中 $P$ 是首一多项式，且 $P'(t)$ 在 $B$ 中可逆. 再设 $T$ 仿射，其环为 $R$，$T_0$ 由幂零理想 $I$ 给出. 在 $R$ 中选取 $u_0(t)$ 的提升 $r_0$. 于是 $P(r_0)\in I$，$P'(r_0)\in R^\times$. 可验证，令 $u(t):=r_0-\frac{P(r_0)}{P'(r_0)}$ 即定义了所需的提升.
反向蕴含（形式平展$\Rightarrow$平展）将由下面的光滑情形得到，因为平展等价于光滑且非分歧.
\emph{光滑情形.} 证明光滑$\Rightarrow$形式光滑. 设 $U$ 是 $X$ 的开集，使 $f|_U$ 是一个平展态射与态射 $\mathbb A_S^n\longrightarrow S$ 的复合. 设 $V_0$ 是 $T_0$ 的仿射开集，且包含于 $u_0^{-1}(U)$. 它也诱导 $T$ 的开集 $V$（因为 $T$ 与 $T_0$ 有相同的拓扑空间）. 此时有图
\begin{align*}
\begin{tikzcd}[ampersand replacement=\&]
V_0\arrow[r,"u_0"]\arrow[d] \& U\arrow[d,"f"]\\
V\arrow[r,"v"']\arrow[ur,dotted,"u"] \& S
\end{tikzcd}
\end{align*}
已经证明平展蕴含形式平展. 因而只须补全下图：
\begin{align*}
\begin{tikzcd}[ampersand replacement=\&]
V_0\arrow[r,"u_0"]\arrow[d] \& \mathbb A_S^n\arrow[d,"f"]\\
V\arrow[r,"v"']\arrow[ur,dotted,"u"] \& S
\end{tikzcd}
\end{align*}
而态射 $u_0$ 是 $v|_{V_0}$ 与一个态射 $V_0\longrightarrow\mathbb A^n$ 的乘积，后者由 $\Gamma(V_0,\mathcal O_{T_0})$ 中的 $n$ 个元素给出. 因而，在 $\Gamma(V,\mathcal O_T)$ 中选取这 $n$ 个元素的提升，就得到态射 $T\longrightarrow\mathbb A^n$，它与 $v$ 的乘积即为所求的 $u$.
用上述 $U$ 覆盖 $X$，便得到 $T$ 的仿射开覆盖 $T=\bigcup_iV_i$，其上均已知存在提升 $u$. 但这里出现一个困难，因为没有任何保证说这些提升确实能够粘合. 由~\eqref{eq:ag-3-4-16-2}，对任意 $i,j$，差 $\delta_{ij}=u_i-u_j$ 属于 $\operatorname{Hom}_{V_i\cap V_j}(u_0^*\Omega_{X/S},\mathcal I)$. 因而 $\delta_{ij}$ 构成覆盖
% 原PDF p65.
$(V_i)_i$ 的一个取值于层 $\mathcal H om_{\mathcal O_{T_0}}(u_0^*\Omega_{X/S},\mathcal I)$ 的 Čech 1-闭上链. 这正是需要 $T$ 仿射的假设之处. 事实上，对任意 Noether 仿射概形 $T=\operatorname{Spec}(R)$ 及任意拟凝聚层 $\mathcal F$，有 $H^i(T,\mathcal F)=0$（$i>0$）.（这是因为内射模所伴随的拟凝聚层是松弛层，从而为 $\Gamma(T,-)$-无上同调层，参见 Hartshorne，III.3.5.）于是 1-闭上链 $(\delta_{ij})_{i,j}$ 可以平凡化，即可写成 $\delta_{ij}=a_i-a_j$，其中 $a_i\in\operatorname{Hom}_{V_i}(u_0^*\Omega_{X/S},\mathcal I)$. 因此，局部提升 $u'_i:=u_i-a_i$ 可以粘合.
还须证明形式光滑蕴含光滑. 为此，只须证明 $\Omega_{X/S}$ 局部自由，并且，给定 $X$ 到 $S$-概形 $Y$ 的一个闭 $S$-嵌入 $g:X\hookrightarrow Y$，其理想记为 $\mathcal J$，则列
\begin{align*}
0\longrightarrow\mathcal J/\mathcal J^2\longrightarrow g^*\Omega_{Y/S}\longrightarrow\Omega_{X/S}\longrightarrow0
\end{align*}
正合. 事实上，这样便可应用定理~\ref{item:ag-3-4-9}的判别准则 iv)，只要选取在 $X$ 上光滑的 $Y$，例如形如 $\mathbb A_S^n$ 的 $Y$.
由前面的讨论，只须考虑仿射情形. 令 $X=\operatorname{Spec}(B)$、$S=\operatorname{Spec}(A)$. 先证明 $\Omega_{B/A}$ 是投射 $B$-模. 从 $B$-模的满态射 $p:M\twoheadrightarrow N$ 出发，须证明态射 $\operatorname{Hom}_B(\Omega_{B/A},M)\longrightarrow\operatorname{Hom}_B(\Omega_{B/A},N)$ 仍为满射. 由命题~\ref{item:ag-3-2-7}，此箭头也可写成 $\operatorname{Der}_A(B,M)\longrightarrow\operatorname{Der}_A(B,N)$. 固定一个 $A$-导子 $\partial:B\longrightarrow N$，并考虑下图：
\begin{align*}
\begin{tikzcd}[ampersand replacement=\&]
B\arrow[r,"\operatorname{Id}\oplus\partial"]\arrow[dr,dotted,"u"] \& B\oplus N\\
A\arrow[u,"f"]\arrow[r,"f\oplus0"'] \& B\oplus M\arrow[u,two heads,"\operatorname{Id}\oplus p"']
\end{tikzcd}
\end{align*}
这里 $B\oplus M$ 表示乘法由 $(b,m)(b',m')=(bb',bm'+b'm)$ 给出的 $B$-代数（特别地，$M$ 是平方为零的理想）. $B\oplus N$ 同样定义. $\partial$ 的导子性质蕴含 $\operatorname{Id}\oplus\partial:B\longrightarrow B\oplus N$ 是 $A$-代数态射. 由形式光滑性，存在图中那样的 $u$. 等式 $(\operatorname{Id}\oplus p)\circ u=(\operatorname{Id}\oplus\partial)$ 蕴含 $u$ 形如 $\operatorname{Id}\oplus\partial'$；计算表明，$\partial':B\longrightarrow M$ 满足 $\partial'(bb')=b\partial'(b')+b'\partial'(b)$. 等式 $u\circ f=f\oplus0$ 蕴含 $\partial'|_A=0$，所以 $\partial'$ 是 $A$-导子，并且满足 $\partial=p\circ\partial'$. 这样便证明了 $\Omega_{B/A}$ 在 $B$ 上投射.
接着证明上述列的正合性. 令 $Y=\operatorname{Spec}(C)$，$I$ 是 $C$ 中使 $B=C/I$ 的理想. 只须给出态射 $\delta:I/I^2\longrightarrow C/I\otimes_C\Omega_{C/A}$ 的一个左逆. 考虑图
\begin{align*}
\begin{tikzcd}[ampersand replacement=\&]
B\arrow[r,equal]\arrow[dr,dotted,"u"] \& C/I\\
A\arrow[u,"f"]\arrow[r] \& C/I^2\arrow[u,two heads]
\end{tikzcd}
\end{align*}
$B$ 在 $A$ 上形式光滑，保证存在图中那样的 $u$. 由此得到分解 $C/I^2=B\oplus I/I^2$，从而得到包含映射
% 原PDF p66.
$I/I^2\subset C/I^2$ 的一个收缩映射 $C/I^2\twoheadrightarrow I/I^2$. 考虑映射 $\partial:C\twoheadrightarrow C/I^2\twoheadrightarrow I/I^2$. 它是 $A$-导子，因而诱导 $C/I$-态射
\begin{align*}
\gamma:C/I\otimes_C\Omega_{C/A}\longrightarrow I/I^2.
\end{align*}
证明 $\gamma\circ\delta=\operatorname{Id}_{I/I^2}$. 由 $\delta$ 的定义，这归结为证明 $\partial$ 在 $I$ 上的限制是典范投影 $I\twoheadrightarrow I/I^2$，而这直接来自构造.
\end{proof}
一般地（不作 Noether 假设），若一个态射形式光滑且有限表现，就称它光滑.
\begin{example}\label{item:ag-3-4-17}
\begin{enumerate}[label=(\roman*)]
\item 若 $X$ 是 $\mathbb Z_p$ 上的光滑概形，则映射 $X(\mathbb Z_p)\longrightarrow X(\mathbb F_p)$ 满射.
\item ``闭点处的纤维''函子，从 $\mathbb Z_p$ 上的平展态射范畴到 $\mathbb F_p$ 上的平展态射范畴，是范畴等价.\unskip\nobreak\hfill\ensuremath{\diamond}
\end{enumerate}
\renewcommand{\envenddiamond}{}
\end{example}